% ============================================================
% P10 — Keranjang Belanja
% ============================================================
\flowmeta{P10}{Keranjang Belanja}{Pelanggan/Tamu}

\begin{flowsection}{Tujuan}
Mengelola daftar produk yang akan dibeli sebelum melanjutkan ke checkout.
\end{flowsection}

\begin{flowsection}{Prasyarat}
Boleh tanpa login (tamu); disarankan login agar keranjang tersimpan pada akun.
Produk tersedia stoknya.
\end{flowsection}

\begin{flowsection}{Langkah Mengelola Keranjang}
\begin{enumerate}
  \item Tambahkan produk dari halaman detail (\textbf{+ Keranjang}).
  \item Buka halaman \texttt{/cart}.
  \item Ubah jumlah dengan tombol \textbf{--} / \textbf{+}; sistem memvalidasi stok.
  \item Centang item (atau \textbf{Pilih Semua}) untuk dihitung ke total.
  \item (Opsional) masukkan \textbf{Kode Kupon} lalu \textbf{Terapkan}.
  \item Klik \textbf{Lanjut ke Pembayaran} untuk menuju checkout (P11).
\end{enumerate}
\end{flowsection}

\shot{gambar/pelanggan/p10-1.png}{Halaman keranjang menampilkan daftar item beserta harga}{fig:p10-1}

\begin{flowsection}{Hasil yang Diharapkan}
Item tampil dengan subtotal; total pembayaran terbarui saat item dipilih; stok
divalidasi secara langsung dari data inventori server.
\end{flowsection}

\shot{gambar/pelanggan/p10-2.png}{Item dipilih sehingga Total Pembayaran terhitung}{fig:p10-2}

\begin{flowsection}{Penanganan Error dan Kasus Khusus}
\begin{itemize}
  \item Stok tidak cukup $\rightarrow$ jumlah dibatasi/pesan peringatan.
  \item Keranjang kosong $\rightarrow$ tampil empty-state ``Keranjang belanjamu kosong''.
  \item Kupon tidak valid $\rightarrow$ muncul pesan kesalahan kupon.
\end{itemize}
\end{flowsection}

\begin{flowsection}{Kaitan Modul}
FE \texttt{CartPage.jsx}, \texttt{cartService.js}. BE \texttt{CartController}. Rute
\texttt{/api/cart}, \texttt{/api/cart/items}, \texttt{/api/vouchers/validate}.
\end{flowsection}

\docver{v1.0}{Sep 2026}
